\chapter{Data Engineering}
\label{ch:data-engineering}

\section{Data Sources}

CAAS ingests data from three independent sources:

\textbf{1. PM2.5 Measurements — PCD Thailand.}
Daily PM2.5 concentration data for Chiang Mai stations 35T and 36T was obtained
from the Pollution Control Department's air4thai portal, covering 2011--2025.
Historical data was provided as Excel files; live daily updates are fetched via
the air4thai JSON API (\texttt{fetch\_pm25\_live.py}).

\textbf{2. Meteorological Data — Open-Meteo.}
Historical and forecast weather variables were obtained from the Open-Meteo API,
an open-access service providing ERA5-quality reanalysis data at daily resolution.
Variables collected include temperature (max/min/mean), humidity (max/min),
precipitation, wind speed, wind direction, surface pressure, cloud cover, and
evapotranspiration (\texttt{fetch\_weather.py}).

\textbf{3. Fire Hotspot Data — NASA FIRMS.}
Daily fire hotspot counts within 50~km and 100~km radii of Chiang Mai were
obtained from the NASA FIRMS API (VIIRS/SNPP satellite product). For each
hotspot radius, mean Fire Radiative Power (FRP) was also extracted as a proxy
for fire intensity. Historical data was batch-fetched by year
(\texttt{run\_firms\_batched.py}); daily updates use \texttt{fetch\_firms.py}.

\section{Data Ingestion Pipeline}

Figure~\ref{fig:daily_pipeline} illustrates the daily ingestion workflow.
For historical backfill, \texttt{run\_pipeline.sh} orchestrates the full sequence
(Steps 1--4) in order. For production operation, a GitHub Actions cron job
(\texttt{daily\_pipeline.yml}) runs every day at 06:00 Bangkok time (23:00 UTC).

The ingestion pipeline:
\begin{enumerate}
    \item Fetches yesterday's PM2.5, weather, and FIRMS data in parallel.
    \item Appends new rows to the consolidated CSVs, deduplicating by date.
    \item Rebuilds the feature matrix (\texttt{build\_features.py}).
    \item Runs inference with the current champion model.
    \item Uploads all results to Amazon S3.
    \item Runs the Evidently drift check and triggers retraining if thresholds are exceeded.
\end{enumerate}

\begin{figure}[H]
    \centering
    \includegraphics[width=0.85\textwidth]{fig_daily_pipeline}
    \caption{CAAS daily data pipeline — from ingestion to drift monitoring.}
    \label{fig:daily_pipeline}
\end{figure}

\section{Feature Engineering}

The feature matrix (\texttt{build\_features.py}) produces 45 features from the
three raw sources, grouped as follows:

\begin{table}[H]
\centering
\caption{CAAS Feature Groups}
\label{tab:features}
\begin{tabular}{llr}
\toprule
\textbf{Group} & \textbf{Features} & \textbf{Count} \\
\midrule
PM2.5 lags          & pm25\_lag1/3/7/14/30                                   & 5  \\
Rolling statistics  & roll3/7/14/30\_mean, roll7/14\_std                     & 6  \\
Temporal encodings  & month, week, day\_of\_year, sin/cos encodings, is\_haze\_season & 8  \\
Weather             & temp (max/min/mean), humidity (max/min), precipitation,
                      wind speed, wind direction (sin/cos), pressure,
                      cloud cover, evapotranspiration                        & 11 \\
Fire (FIRMS)        & hotspot\_50km, hotspot\_100km, mean\_frp\_50km,
                      hotspot\_7d\_roll, hotspot\_14d\_roll, fire\_flag      & 6  \\
Alert lags          & alert\_lag1/3/7                                         & 3  \\
Cross features      & lag1$\times$temp, lag1$\times$wind, roll7$\times$fire,
                      precip\_flag, humidity$\times$wind                     & 5  \\
\midrule
\textbf{Total}      &                                                         & \textbf{45} \\
\bottomrule
\end{tabular}
\end{table}

Key engineering decisions:
\begin{itemize}
    \item \textbf{No data leakage:} Lag and rolling features use only past values.
          Target columns (pm25\_t1, pm25\_t3, pm25\_t7) are forward-shifted.
    \item \textbf{Fire accumulation features:} 7-day and 14-day rolling hotspot
          counts capture the cumulative smoke load better than a single-day count.
    \item \textbf{Cyclic temporal encodings:} Month and day-of-year are encoded
          as sine/cosine pairs to preserve calendar continuity (December-January
          boundary).
    \item \textbf{Cross features:} lag1$\times$temp and roll7$\times$fire capture
          interactive effects between PM2.5 persistence and environmental conditions.
\end{itemize}

\section{Data Storage and Versioning}

All data artefacts are stored in Amazon S3 with a structured prefix layout:

\begin{verbatim}
s3://caas-mlops-bucket/
    data/raw/         ← original Excel and CSV files
    data/processed/   ← pm25_consolidated.csv, weather, FIRMS
    data/features/    ← features.csv (daily snapshot)
    models/           ← xgboost_t1/3/7.json (champion models)
    results/          ← forecasts, drift reports, summaries
    mlflow/           ← MLflow artefacts
\end{verbatim}

S3 bucket versioning is enabled, preserving every version of model files
and allowing rollback to any previous champion.

\section{Pipeline Orchestration}

Pipeline scheduling is handled by GitHub Actions:
\begin{itemize}
    \item \textbf{daily\_pipeline.yml} --- cron 06:00 Bangkok time, runs ingest +
          inference + upload + drift check.
    \item \textbf{retrain.yml} --- triggered annually (December~1) or on-demand
          via \texttt{repository\_dispatch} from the drift monitor.
    \item \textbf{test.yml} --- runs pytest suite on every push to main.
\end{itemize}

This architecture was chosen over Apache Airflow for simplicity within the
capstone scope. Airflow would be appropriate for multi-team production use
requiring finer DAG-level observability and retries.
